\section{Theoretical Foundations}
\label{appendix:theory}

% 定义理论中的谓词记号
\newcommand{\passmark}{\mathcal{P}_{\text{all}}}
\newcommand{\rejectmark}{\mathcal{R}_{\text{all}}}

We formalize the behavior of a Large Language Model (LLM) acting as a theorem prover within the context of specification verification. This abstraction enables rigorous analysis of the fundamental limitations inherent in using LLMs for formal verification tasks.

\subsection{LLM as a Probabilistic Prover}
\label{appendix:llm-model}

Consider a specification $S$ defined as a logical predicate over input-output pairs. Given a test case $(i, o)$ consisting of an input $i$ and output $o$, we say that $S$ \emph{accepts} the test case, denoted $S(i, o) = \top$, if the input-output pair satisfies the specification. Conversely, $S$ \emph{rejects} the test case, denoted $S(i, o) = \bot$, if the pair violates the specification.

We model an LLM-based prover as a probabilistic program $\mathcal{M}$ that takes as input a specification $S$ and a test case $(i, o)$, and attempts to produce a proof of acceptance.

\begin{definition}[LLM Proving Oracle]
\label{def:llm-oracle}
Let $\mathcal{M}$ be an LLM-based prover. For a specification $S$ and test case $(i, o)$, define the random variable $\textsc{Verify}_{\mathcal{M}}(S, i, o) \in \{\texttt{proof}, \bot\}$ as follows:
\begin{itemize}[leftmargin=1.5em, itemsep=4pt, topsep=4pt]
    \item If $S(i, o) = \top$ (true acceptance):
    \begin{align}
    \Pr[\textsc{Verify}_{\mathcal{M}}(S, i, o) = \texttt{proof}] &= p_{S,i,o} \\
    \Pr[\textsc{Verify}_{\mathcal{M}}(S, i, o) = \bot] &= 1 - p_{S,i,o}
    \end{align}
    where $p_{S,i,o} \in [0, 1]$ is the \emph{proof success probability}, depending on $S$, $(i, o)$, and $\mathcal{M}$.
    \item If $S(i, o) = \bot$ (true rejection):
    \begin{align}
    \Pr[\textsc{Verify}_{\mathcal{M}}(S, i, o) = \texttt{proof}] &= 0 \\
    \Pr[\textsc{Verify}_{\mathcal{M}}(S, i, o) = \bot] &= 1
    \end{align}
\end{itemize}
\end{definition}

The second case captures the fundamental \emph{soundness} property: the LLM cannot produce a valid proof for a false proposition, as any generated proof would be rejected by the Rocq type checker.

\subsection{The Undecidability of Verification Failure}

A critical consequence of our model is that verification failure is fundamentally uninformative about specification correctness.

\begin{theorem}[Verification Undecidability]
\label{thm:undecidability}
When $\textsc{Verify}_{\mathcal{M}}(S, i, o) = \bot$, it is impossible to determine whether $S(i, o) = \top$ or $S(i, o) = \bot$ without additional information about $p_{S,i,o}$.
\end{theorem}

\begin{proof}
Consider two scenarios producing identical observable outcomes:
\begin{itemize}[leftmargin=1.5em, itemsep=2pt, topsep=2pt]
    \item $S(i, o) = \top$ but $p_{S,i,o}$ is small, i.e., the proof is difficult.
    \item $S(i, o) = \bot$, i.e., no valid proof exists.
\end{itemize}
Both yield $\textsc{Verify}_{\mathcal{M}}(S, i, o) = \bot$. Without knowledge of $p_{S,i,o}$, these are observationally indistinguishable.
\end{proof}

This theorem reveals a fundamental Finding \ref{duality} \emph{evaluation duality}: when verification fails, the fault may lie in either the specification (generator error) or the prover (prover limitation). This provides the formal basis for the ``Evaluation Duality'' observation in Section~\ref{sec:rq1}.

A key insight is the \emph{information asymmetry} between outcomes:
\begin{itemize}[leftmargin=1.5em, itemsep=2pt, topsep=2pt]
    \item Positive evidence ($\textsc{Verify} = \texttt{proof}$) guarantees $S(i,o) = \top$ by Rocq's soundness.
    \item Negative evidence ($\textsc{Verify} = \bot$) provides no definitive information.
\end{itemize}
This asymmetry motivates using successful proof construction as the primary evaluation criterion.

\subsection{Factors Affecting Proof Success Probability}

The proof success probability $p_{S,i,o}$ varies based on multiple factors:
\begin{itemize}[leftmargin=1.5em, itemsep=4pt, topsep=4pt]
    \item \emph{Proof complexity}: $p_{S,i,o}$ decreases as proof difficulty increases, including proof length, need for auxiliary lemmas, and sophisticated tactics. Our experiments in Section~\ref{sec:rq2} support this: all models show substantial drops from \textsc{Pass}$_{\text
    {first}}$ to \textsc{Pass}$_{\text{all}}$ (e.g., Gemini 3 Pro Preview: 60.98\% $\to$ 28.05\%).
    \item \emph{Specification style}: executable components (\texttt{Fixpoint}) often admit simpler proofs via computation, increasing $p_{S,i,o}$; pure relational specifications (\texttt{Inductive}) may require explicit inductive reasoning, decreasing $p_{S,i,o}$. This is empirically validated in Section~\ref{sec:rq4}.
    \item \emph{Test case complexity}: larger inputs or edge cases require more intricate proof steps, reducing $p_{S,i,o}$.
    \item \emph{Model capability}: different LLMs exhibit varying $p_{S,i,o}$ for the same task, as shown in Table~\ref{tab:verification}.
\end{itemize}

\subsection{Formal Analysis of the \ours{} Framework}

We provide formal justification for our evaluation stages (Section~\ref{sec:framework}). To distinguish \emph{metrics} (percentages in experiments) from \emph{predicates} (boolean functions), we use:
\begin{itemize}[leftmargin=1.5em, itemsep=2pt, topsep=2pt]
    \item $\passmark(S), \rejectmark(S) \in \{\top, \bot\}$: predicates on specification $S$.
    \item \textsc{Pass}$_{\text{all}}$, \textsc{Reject}$_{\text{all}}$: metrics representing the proportion of specifications satisfying the predicates.
\end{itemize}

\paragraph{The $\passmark$ predicate.}
Let $\mathcal{T}^{+} = \{(i_1, o_1), \ldots, (i_n, o_n)\}$ be the positive test cases. Define:
\begin{equation}
\passmark(S) = \top \iff \forall (i, o) \in \mathcal{T}^{+} : \textsc{Verify}_{\mathcal{M}}(S, i, o) = \texttt{proof}
\end{equation}
The metric \textsc{Pass}$_{\text{all}} = |\{S : \passmark(S) = \top\}| / N$ where $N$ is the total number of specifications. By Rocq's soundness, $\passmark(S) = \top$ implies $S$ accepts all positive test cases: $\forall (i, o) \in \mathcal{T}^{+} : S(i, o) = \top$. Note that $\passmark(S) = \bot$ does not imply incorrectness---it may indicate low $p_{S,i,o}$ for some test cases.

\paragraph{The $\rejectmark$ predicate.}
Let $\mathcal{T}^{-} = \{(i'_1, o'_1), \ldots, (i'_m, o'_m)\}$ be the negative cases (Section~\ref{sec:neg-cases}). Define:
\begin{equation}
\rejectmark(S) = \top \iff \forall (i', o') \in \mathcal{T}^{-} : \textsc{Verify}_{\mathcal{M}}(S, i', o') = \bot
\end{equation}
If $\textsc{Verify}_{\mathcal{M}}(S, i', o') = \texttt{proof}$ for any $(i', o') \in \mathcal{T}^{-}$, then $S$ is provably overly permissive. Conversely, $\rejectmark(S) = \top$ does not guarantee $S$ rejects all negative cases---the LLM may fail to find proofs even when they exist.

\paragraph{Candidate specifications.}
A specification $S$ with $\passmark(S) = \rejectmark(S) = \top$ is a \emph{candidate specification} (corresponding to specifications passing \textsc{Reject}$_\text{all}$ in Table~\ref{tab:generation}). Such specifications:
\begin{itemize}[leftmargin=1.5em, itemsep=2pt, topsep=2pt]
    \item provably accept all positive test cases;
    \item are not provably overly permissive;
    \item may still be incorrect in ways not captured by the test suite.
\end{itemize}

\subsection{Why Full Equivalence Proving is Impractical}

Consider proving equivalence $\forall\, i, o : S_1(i, o) \Leftrightarrow S_2(i, o)$. Unlike test-case verification, this requires handling arbitrary inputs through induction. Let $p_{\text{equiv}}$ denote the success probability; empirically $p_{\text{equiv}} \ll p_{S,i,o}$.

Our results in Appendix~\ref{appendix:human-llm} confirm this: among specifications passing all test-based criteria, only 31.0\% yield provably isomorphic specifications. This illustrates the
Finding \ref{trade-off} \emph{verification complexity trade-off}: as verification complexity increases, proof success probability decreases, causing systematic underestimation of specification quality.

\subsection{Design Principles}

Our analysis yields four design principles:

\textbf{Exploit asymmetry.} Maximize use of positive evidence while acknowledging uncertainty in negative evidence.

\textbf{Calibrate difficulty.} Avoid verification tasks that are too difficult; overly complex proof obligations lead to uniformly low $p_{S,i,o}$, preventing meaningful differentiation between specifications.

\textbf{Use strongest prover.} Maximize $p_{S,i,o}$ to minimize false negatives from prover limitations.

\textbf{Anchor the duality.} Use known-good artifacts (human specifications) to anchor assessment.

The \ours{} framework embodies these principles through test-case-based verification, Gemini 3 Pro Preview as prover (Section~\ref{sec:rq1}), human specification anchors (Section~\ref{sec:rq3}), and epistemic humility about failed proofs.
